\chapter*{Executive Summary}
\addcontentsline{toc}{chapter}{Executive Summary}
\textbf{Problem.} Store managers cannot easily tell whether customers on the floor are being served, how staff time is split across the floor, or which aisles are left uncovered. CCTV footage holds this information but is watched only after incidents, and commercial retail-analytics products need extra hardware or budgets that small stores do not have.

\textbf{Solution.} RetailPulse reuses existing cameras. A computer-vision pipeline detects and tracks people, labels them as customers or employees from appearance cues, and confirms interactions with a pose check. Results are written to a PostgreSQL event store. On top of the store sit a Customer Behavior agent and an Employee Activity agent, an orchestrator that routes manager questions and checks every number against its source, a signal policy that raises unattended-customer and unstaffed-zone events, and a dispatch service that assigns an eligible employee directly. A Next.js web application gives managers a dashboard and employees a task workspace.

\textbf{Progress since the mid evaluation.} The mid evaluation found the figures and the software requirements specification incorrect. Both were rebuilt for this report: the requirements specification now has sequentially numbered, atomic and testable requirements with measurable non-functional targets (Chapter~\ref{ch:srs}), and every design figure was redrawn from the current code (Chapter~\ref{ch:design}). In the same period the team moved the system from SQLite to PostgreSQL, added the event store, the agents and their grounding checks, direct dispatch with duplicate protection, the four live occupancy categories, map registration, security fixes and continuous integration, and evaluated the vision pipeline on 25 labelled clips.

\textbf{Results.} On held-out clips the interaction detector scores F1 0.450 with role classification correct for all 56 matched tracks; it runs at 0.60 times real time on a CPU. All three continuous-integration jobs pass. A rehearsal of the ten acceptance scenarios on a synthetic store day passes eight; the remaining two need a store map with camera calibration and two human reviewers.

\textbf{Remaining work.} Acceptance on real store footage, a deployed backend, the store floor plan, a second label annotator and the model freeze are planned for FYP-2 (Chapter~\ref{ch:conclusion}).
